% Datos de portada, cabeceras y metadatos del Tema 5 (primera parte)
\renewcommand{\theproblema}{5.\arabic{problema}}
\newcommand{\PortadaAsignatura}{Electromagnetismo I}
\newcommand{\PortadaTema}{Tema 5 \;$\cdot$\; Parte 1}
\newcommand{\PortadaTitulo}{Energía\\[2pt]y fuerzas}
\newcommand{\PortadaSubtitulo}{Energía de cargas y distribuciones, energía del campo, sistemas de conductores\\
  y fuerzas por el método de los trabajos virtuales}
\newcommand{\PortadaImagen}{t5-portada-oscuro}
\newcommand{\PortadaFormulas}{U_e=\frac12\sum_{i}q_i\phi_i,\qquad U_e=\frac{\eo}{2}\int E^2\dd V,\qquad
  \vF_e=-\grad U_e\big|_{Q},\qquad \vF_e=+\grad U_e\big|_{\phi}}
\newcommand{\PortadaAutor}{Daniel Martín Guerrero}
\newcommand{\PortadaProfesor}{Luis López Díaz}
\newcommand{\PortadaTipo}{Apuntes formalizados}
\newcommand{\PortadaCurso}{Salamanca, curso 2026/2027}
\renewcommand{\cabizq}{Tema 5 $\cdot$ Energía y fuerzas (I)}
\renewcommand{\cabder}{Electromagnetismo I}
\hypersetup{pdftitle={Tema 5. Energía y fuerzas (primera parte)},
  pdfauthor={Daniel Martín Guerrero},
  pdfsubject={Electromagnetismo I. Apuntes formalizados con seminarios prácticos},
  pdfkeywords={energía electrostática, densidad de energía, conductores, trabajos virtuales, Earnshaw}}
% Notación propia del tema
\newcommand{\Wext}{W_{\mathrm{ext}}}
\newcommand{\Wmec}{W_{\mathrm{mec}}}
\newcommand{\Wbat}{W_{\mathrm{bat}}}
\newcommand{\Uint}{U_{\mathrm{int}}}
\newcommand{\vFe}{\vF_e}
